\section{The direct Gauss--Ramanujan route}
\label{sec:gauss-complementarity}

The norm route does not use the conductor sensitivity of the additive
Gauss factor.  We now exploit it before taking a fiber norm.  For
\(\chi\in\mathfrak X(q)\), define
\begin{equation}
 G_q(\chi;t)=\sum_{\xi\in G_q}\chi(\xi)e_q(t\xi).
 \label{eq:generalized-Gauss}
\end{equation}

\begin{lemma}[Induced Gauss factor at every frequency]
\label{lem:induced-Gauss}
Let \(q\) be squarefree, let \(f=\cond\chi\mid q\), let
\(\chi^*\bmod f\) induce \(\chi\), and put \(d=q/f\).  Then
\begin{equation}
 G_q(\chi;t)=
 \begin{cases}
  \overline{\chi^*(t)}\chi^*(d)\tau_f(\chi^*)c_d(t),
     &(t,f)=1,\\
  0,&(t,f)>1.
 \end{cases}
 \label{eq:induced-Gauss}
\end{equation}
Here \(\tau_1(\mathbf1)=1\).  Consequently, for \(f>1\),
\begin{equation}
 |G_q(\chi;t)|
 =\sqrt f\,|c_d(t)|\one_{(t,f)=1}.
 \label{eq:induced-Gauss-size}
\end{equation}
\end{lemma}

\begin{proof}
CRT factors \eqref{eq:generalized-Gauss} into a primitive Gauss sum
modulo \(f\) and a Ramanujan sum modulo \(d\).  The first factor
vanishes when \((t,f)>1\); otherwise the usual change of variables
gives \(\overline{\chi^*(t)}\chi^*(d)\tau_f(\chi^*)\).
The second factor is \(c_d(t)\).  Primitive Gauss magnitude is
\(\sqrt f\).
\end{proof}

Since \((H,q)=1\), replacing \(t\) by \(r\overline H\) does not change
the condition \((r,f)=1\), and it changes only a unit-modulus character
phase.

\begin{lemma}[Ramanujan-weighted nonzero frequencies]
\label{lem:Ramanujan-Fourier}
If \(s\) is squarefree, \(K>0\), and \(F\) is fixed and smooth with
rapidly decaying Fourier transform, then
\begin{equation}
 \boxed{
 \sum_{r\ne0}|c_s(r)|\,|\widehat F(r/K)|
 \ll_F K\tau(s).}
 \label{eq:Ramanujan-Fourier}
\end{equation}
The exclusion of \(r=0\) is essential.
\end{lemma}

\begin{proof}
For squarefree \(s\),
\begin{equation}
 |c_s(r)|=\varphi((s,r))
 \le(s,r)=\sum_{e\mid(s,r)}\varphi(e).
 \label{eq:Ramanujan-gcd}
\end{equation}
For every \(e\ge1\), rapid decay gives uniformly in \(K>0\)
\begin{equation}
 \sum_{\substack{r\ne0\\e\mid r}}|\widehat F(r/K)|
 =\sum_{k\ne0}|\widehat F(ke/K)|
 \ll_F \frac Ke.
 \label{eq:Fourier-multiples}
\end{equation}
When \(e>K\), take two or more powers of decay; this is precisely where
removing \(r=0\) matters.  Insert \eqref{eq:Ramanujan-gcd}, interchange
the divisor and frequency sums, and use
\(\sum_{e\mid s}\varphi(e)/e\le\tau(s)\).
\end{proof}

\subsection{One exact pair--conductor block}

Fix \(g,a,b\), let \(u=ga\), \(v=gb\), \(q=gab\), and fix exact
coordinate conductors
\begin{equation}
 f_1=\cond\alpha,\qquad f_2=\cond\beta,\qquad
 f=f_1f_2,\qquad d=\frac qf.
 \label{eq:exact-coordinate-conductors}
\end{equation}
Define
\begin{equation}
 A_{g,a,b;f_1,f_2}^2
 :=|\vartheta_R(u)\vartheta_R(v)|^2
 \sum_{\substack{\alpha\in\mathfrak X(u),\ \cond\alpha=f_1\\
                  \beta\in\mathfrak X(b),\ \cond\beta=f_2}}
 |C_{g,a,b}(\alpha,\beta)|^2.
 \label{eq:exact-block-amplitude}
\end{equation}
By \eqref{eq:conductor-product}, all characters in this block have
global conductor \(f\).  Let
\(D_{g,a,b;f_1,f_2}\) be its inverse multiplicative-character transform,
including the factor \(\vartheta_R(u)\vartheta_R(v)\) and the
unit-modulus phase \(\overline{\chi(-h)}\), and define
\begin{equation}
 S_{g,a,b;f_1,f_2}(r)
 :=\sum_{\xi\in\cU_q}D_{g,a,b;f_1,f_2}(\xi)
 e_q(r\overline H\xi).
 \label{eq:exact-block-additive-transform}
\end{equation}
Multiplicative inversion and \cref{lem:induced-Gauss} show that
\begin{equation}
 |S_{g,a,b;f_1,f_2}(r)|
 \le \frac f{\varphi(q)}|c_d(r)|
 A_{g,a,b;f_1,f_2}.
 \label{eq:exact-block-S}
\end{equation}
Indeed, the number of primitive characters of conductor \(f\) is at
most \(f\); Cauchy--Schwarz combines its square root with the Gauss
magnitude \(\sqrt f\) to give exactly \(f\).  Terms with \((r,f)>1\)
vanish and may be retained harmlessly on the right.

Insert \eqref{eq:exact-block-S} into \eqref{eq:scalar-packet}.  Since
\((H,d)=1\), multiplicativity gives
\(|c_H(r)c_d(r)|=|c_{Hd}(r)|\).  Apply
\cref{lem:Ramanujan-Fourier} with \(s=Hd\) and \(K=Hq/J\).  The factors
\(J/(Hq)\) and \(Hq/J\) cancel, yielding
\begin{align}
 |\cL_{g,a,b;f_1,f_2}|
 &\ll_{F,H}\frac f{\varphi(q)}\tau(Hd)
 A_{g,a,b;f_1,f_2}\notag\\
 &\ll_{F,H,\eps}X^\eps \frac fq
 A_{g,a,b;f_1,f_2}.
 \label{eq:exact-direct-bound}
\end{align}
The second line uses \(q/\varphi(q)\ll_\eps X^\eps\).  This estimate
uses all nonzero frequencies, not merely those coprime to \(q\).

\subsection{Dyadic conductor aggregation}

Exact conductors must be split before dyadic aggregation: the
Ramanujan factor \(c_{q/(f_1f_2)}(r)\) changes with the exact pair
\((f_1,f_2)\).  The number of exact conductor pairs and lcm divisor
pairs associated with a fixed squarefree \(q\) is
\(\tau(q)^{O(1)}\), hence \(X^{o(1)}\).

Sum \eqref{eq:exact-direct-bound} over
\(f_i\asymp F_i\), over divisor pairs, and over \(q\asymp Y\).
Weighted Cauchy--Schwarz gives
\begin{align}
 \sum \frac{A_{g,a,b;f_1,f_2}}q
 &\le
 \left(\sum \frac{A_{g,a,b;f_1,f_2}^2}{\varphi(q)}\right)^{1/2}
 \left(\sum \frac{\varphi(q)}{q^2}\right)^{1/2}\notag\\
 &\ll_\eps X^\eps\cE_Y(F_1,F_2)^{1/2}.
 \label{eq:weighted-q-Cauchy}
\end{align}
More explicitly, if \(I=(g,a,b,f_1,f_2)\) is the full summation index,
then
\[
 \sum_I\frac{\varphi(q_I)}{q_I^2}
 \ll X^\eps\sum_{q\asymp Y}\frac{\varphi(q)}{q^2}
 \ll X^\eps.
\]
There is no factor \(\sqrt Y\).

\begin{theorem}[Direct-Gauss conductor bound]
\label{thm:direct-Gauss-bound}
Every nonprincipal dyadic coordinate-conductor block satisfies
\begin{equation}
 \boxed{
 \mathfrak L_Y(F_1,F_2)
 \ll_\eps X^\eps QF_1F_2
 \sqrt{\cK_Q(F_1)\cK_Q(F_2)}.}
 \label{eq:direct-Gauss-bound}
\end{equation}
\end{theorem}

\begin{proof}
Equations \eqref{eq:exact-direct-bound} and
\eqref{eq:weighted-q-Cauchy} give
\[
 \mathfrak L_Y(F_1,F_2)
 \ll_\eps X^\eps F_1F_2\cE_Y(F_1,F_2)^{1/2}.
\]
Insert \eqref{eq:tensor-energy-bound}.
\end{proof}

Combining \eqref{eq:tensor-norm-scalar} and
\eqref{eq:direct-Gauss-bound} for the same object gives the compact
cell estimate
\begin{equation}
 \boxed{
 \mathfrak L_Y(F_1,F_2)
 \ll_\eps X^\eps Q
 \sqrt{\cK_Q(F_1)\cK_Q(F_2)}
 \min\{\sqrt{JY},F_1F_2\}.}
 \label{eq:combined-cell-bound}
\end{equation}

The block \((F_1,F_2)=(1,1)\) is the global principal character and
belongs to the mean already handled in \eqref{eq:imported-mean-bound}.
The one-sided blocks \((1,F_2)\) and \((F_1,1)\) with the other
conductor nontrivial remain in \cref{thm:direct-Gauss-bound}.
